\subsection{The linearisation of a graded bundle}
 In this subsection we shall recall the construction of the \emph{linearisation functor} given in~\cite{Bruce:2014}. In order to avoid some notational clashes, we have changed some of the notation as compared with previous works.

The tangent bundle $\sT F_{k}$ of a graded bundle naturally has the structure of a double graded bundle. The f\/irst weight vector f\/ield is simply the \emph{tangent lift} \cite{Grabowski:1995,Yano:1973} $\Delta^{1} := \rmd_{\sT}\Delta_{F_{k}}$ of the weight vector f\/ield $\Delta_{F_{k}}$ on $F_{k}$, and the second being the natural Euler vector f\/ield on the tangent bundle $\Delta^{2} := \Delta_{\sT F_{k}}$. The tangent bundle of a~graded bundle is a canonical example of a graded-linear bundle, i.e., a graded bundle equipped with a compatible linear structure on the total space.

\begin{Definition}[\cite{Bruce:2014}] A double graded bundle $\sF = (F; \Delta^{1}, \Delta^{2})$ such that~$\zD^1$ is of degree $k-1$ and $\zD^2$ is of degree~1, i.e., $\Delta^{2}$ is an Euler vector f\/ield, will be referred to as a \emph{graded-linear bundle}, or for short a $\GrL$-bundle.
\end{Definition}

It follows that $\sF$ is of total weight $\le k$ with respect to $\Delta := \Delta^{1} + \Delta^{2}$ and a vector bundle structure
\begin{gather*}
\textnormal{p}^{F}_{[\Delta^{2} \leq 0]}\colon \  F \rightarrow F\big[\Delta^{2} \leq 0\big].
\end{gather*}
with respect to the projection onto the submanifold $F[\Delta^{2} \leq 0]$ which inherits a graded bundle structure of degree $k-1$.

\begin{Example} Consider the vertical bundle with respect to the projection $\tau\colon F_{k} \rightarrow M$. The weight vector f\/ields on the vertical bundle $\sV F_{k} \subset \sT F_{k}$ are simply the appropriate restrictions of those on the tangent bundle. \emph{Via} passing to the total weight, we can view~$\sV F_{k}$ as a graded bundle of degree $k+1$. However, it will be useful to shift the f\/irst component of bi-weight to allow us to consider the vertical bundle as a graded bundle of degree $k$. We will always consider this shifted weight when encountering a vertical bundle of a~graded bundle. In other words, if we take $\Delta^1\in \XX(\sT F_k)$ the tangent lift of the weight vector f\/ield and $\Delta^2\in \XX(\sT F_k)$ the Euler vector f\/ield associated with the tangent bundle structure on $F_k$ we get
\begin{gather*}
\sV F_k = \sT F_k\big[\Delta^1-\Delta^2 \geq 0\big],
\end{gather*}
equipped with the weight vector f\/ield $\Delta^1_{\sV F_k}:= (\Delta^1-\Delta^2)|_{\sV F_k}$ of degree $k$ and the Euler vector f\/ield $\Delta^2|_{\sV F_k}$. It is crucial here that, although it is not a combination with non-negative coef\/f\/icients, $\Delta^1-\Delta^2$ is still a weight vector f\/ield on $\sV F_k$.
\end{Example}

\begin{Example}
 Consider a degree two graded bundle $F_{2}$ which we equip with homogeneous local coordinates $(x,y,z)$ of weight $0$, $1$ and $2$, respectively. Taking the tangent functor combined with the necessary shift in the weight, produces homogeneous local coordinates
 \begin{gather*}
 \Big(\underbrace{x}_{(0,0)}, \hspace{5pt} \underbrace{y}_{(1,0)}, \hspace{5pt} \underbrace{z}_{(2,0)}; \hspace{5pt} \underbrace{\dot{x}}_{(-1,1)},\hspace{5pt}\underbrace{\dot{y}}_{(0,1)}, \hspace{5pt} \underbrace{\dot{z}}_{(1,1)}\Big).
 \end{gather*}
The removal of the coordinate $\dot{x}$ is well def\/ined (see Def\/inition~\ref{def:removal_coordinates2}). Doing so produces a double graded bundle which we recognize as the standard vertical bundle of $F_{2}$.
 \end{Example}

The \emph{linearisation functor} takes a graded bundle and produces a $\GrL$-bundle. The basic idea is to mimic the canonical embedding
\begin{gather*}
\sT^{k}M \hookrightarrow \sT\big(\sT^{k-1}M\big),
 \end{gather*}
which sends a point in $\sT^k M$ represented by a curve $t \mapsto \gamma(t)$ to a vector tangent to $\sT^{k-1} M$ represented by the curve $s\mapsto [t\mapsto \gamma(s+t)]_{k-1}\in \sT^{k-1} M$. The linearisation of a~graded bundle~$F_{k}$ is then viewed as a reduction of the tangent bundle~$\sT F_{k}$.
\begin{Definition}[\cite{Bruce:2014}]\label{def:linearisation}
The \emph{linearisation of a graded bundle} $F_{k}$ is the $\GrL$-bundle def\/ined as
\begin{gather*}
\pLinr (F_{k}) := \sV F_{k}\big[\Delta_{\sV F_{k}}^{1} \leq k-1\big],
\end{gather*}
so that we have the natural projection $\textnormal{p}^{\sV F_{k}}_{\pLinr(F_{k})}\colon \sV F_{k} \rightarrow \pLinr(F_{k})$. With a small abuse of notation, we would write $\pLinr (F_k) = \sT F_k[0\leq \Delta^1 - \Delta^2 \leq k-1]$, i.e., we get rid of coordinates on~$\sT F_k$ of negative weights and of weight~$k$ with respect to the vector f\/ield $\Delta^1 - \Delta^2$.
\end{Definition}

\begin{Theorem}[\cite{Bruce:2014}]\label{theom:functorial linearisation}
Linearisation is a functor from the category of graded bundles to the category of double graded bundles.
\end{Theorem}

A full proof of the above theorem can be found in \cite{Bruce:2014}. Essentially, the functorial properties of the linearisation follow from the fact that it is constructed from the tangent functor.

Let us brief\/ly describe the local structure of the linearisation. Consider $F_{k}$ equipped with local coordinates $\big(x^{A}, y_{w}^{a}, z^{i}_{k}\big)$, where the weights are assigned as $\w\big(x^A\big)=0$, $\w(y^a_{w}) = w$ $(1\leq w < k)$ and $\w(z^i_k) =k$. It will be convenient to single out the highest weight coordinates, as well as the zero weight coordinates, and so we change our notation slightly. In any natural homogeneous system of coordinates on the vertical bundle $\sV F_{k}$, one projects out the highest weight coordinates on~$F_{k}$ to obtain~$\pLinr(F_{k})$. Thus, on $\pLinr(F_{k})$ we have local homogeneous coordinates
\begin{gather*}
\Big(\underbrace{x^{A}}_{(0,0)}, \hspace{5pt} \underbrace{y^{a}_{w}}_{(w,0)}; \hspace{5pt} \underbrace{\dot{y}^{b}_{w}}_{(w-1,1)}, \hspace{5pt} \underbrace{\dot{z}_{k}^{i}}_{(k-1,1)}\Big).
 \end{gather*}
Note that with this assignment of the weights, the linearisation of a graded bundle of degree~$k$ is itself a graded bundle of degree~$k$ when passing from the bi-weight to the total weight. It is important to note that the linearisation of a graded bundle has the structure of a vector bundle $\pLinr(F_{k}) \rightarrow F_{k-1}$, hence the nomenclature ``linearisation''. The vector bundle structure is clear from the construction by examining the bi-weight. The second leg of $\GrL$-bundle $\pLinr(F_{k})$ is a degree $k-1$ graded bundle with base~$F_1$:
\begin{gather*}
\xymatrix{
\pLinr(F_{k}) \ar[rr]\ar[d] && F_1 \ar[d] \\
F_{k-1} \ar[rr] && M.
}
\end{gather*}
The basic properties of the linearisation functor are summarized in the following.
\begin{Proposition}\label{prop:pLin}\quad
\begin{enumerate}[$(a)$]\itemsep=0pt
\item \label{prop_pL:item_pullback} There exist a canonical vector bundle morphism $p=\textnormal{p}^{\sV F_{k}}_{\pLinr(F_{k})} \colon \sV F_{k} \rightarrow \pLinr(F_{k})$ covering the projection $\tau^k_{k-1}\colon F_k\rightarrow F_{k-1}$ such that $\sV F_{k}$ is the pullback of the vector bundle $\pLinr(F_{k})\rightarrow F_{k-1}$ by $p$:
 \begin{gather*}
\xymatrix{
p^* \pLinr(F_{k}) \simeq \sV F_{k}\ar[rr]^p\ar[d] && \pLinr(F_{k}) \ar[d] \\
F_k \ar[rr]^{\tau^k_{k-1}} && F_{k-1}.
}
\end{gather*}
\item \label{prop_pL:itemTkM} The linearisation $\pLinr(\sT^k M)$ of the $k$-th order tangent bundle of $M$ is isomorphic with $\sT \sT^{k-1} M$ as $\GrL$-bundle. In particular, there is a canonical epimorphism $ p\colon \sV \sT^k M = \ker (\sT \tau^k_M\colon \sT\sT^k M $ $\rightarrow \sT M) \twoheadrightarrow \sT\sT^{k-1} M$. In coordinates,
 \begin{gather*}
 p^*\big(x^{a, (\alpha)}\big) = x^{a, (\alpha)}, \qquad p^*\big(dx^{a, (\beta-1)}\big) = dx^{a, (\beta)},
 \end{gather*}
 where $0\leq \alpha\leq k-1$, $1\leq \beta \leq k$.
\item \label{prop_pL:item embedding} There is a canonical $($total$)$ weight preserving embedding
\begin{gather*}
\iota \colon \ F_{k} \hookrightarrow \pLinr(F_{k})
\end{gather*}
given by the composition of the weight vector field $\Delta_{F_k} \in \Vect(F_k)$ considered as a section of~$\sV F_{k}$ with the natural projection $p\colon \sV F_{k} \rightarrow \pLinr(F_{k})$. In natural local coordinates, the nontrivial part of the embedding is given by
\begin{gather*}
\iota^{*}(\dot{y}^{a}_{w}) = w   y^{a}_{w}, \qquad \iota^{*}(\dot{z}_{k}^{i}) = k  z_{k}^{i}.
\end{gather*}
\item \label{prop_pL:item_pl_tau} The induced morphism $\pLinr(\tau^k_{k-r})\colon \pLinr(F_k)\rightarrow \pLinr(F_{k-r})$ coincides with $\pLinr(F_k) \rightarrow \pLinr(F_k)[\Delta^1_{\pLinr(F_{k})} \leq k-r-1]$, where $\tau^k_{k-r} = p^{F_k}_{[\Delta_{F_k}\leq k-r]}\colon  F_k\to F_{k-r}$ is the canonical projection on a lower graded bundle.
\item \label{prop_pL:item_injective} Given an injective graded bundle morphism $\phi\colon F_k\rightarrow F_k'$, the associated morphism $\pLinr(\phi)$ is injective, as well.
\end{enumerate}
\end{Proposition}

\begin{proof} We shall only give a geometrical explanation of~\eqref{prop_pL:itemTkM}. Another proof can be found in the arXiv preprint version of~\cite{Bruce:2014b}. The statement~\eqref{prop_pL:item_pullback} follows directly from the def\/inition of the linearisation functor as both vector bundles in interest are of the same rank. The statements~\eqref{prop_pL:item_pl_tau} and~\eqref{prop_pL:item_injective} are very easy to proof directly, while~\eqref{prop_pL:item embedding} was carefully proved in~\cite{Bruce:2014}.

Recall~\cite{Leon:1992, Eliopoulos:1966}, a $k$-th tangent bundle $\sT^k M$ is equipped with an endomorphism~$J$ of its tangent bundle, called the canonical \emph{almost tangent structure of order~$k$}, which has the form
\begin{gather*}
J \pa_{x^{A, (\alpha)}} = \pa_{x^{A, (\alpha+1)}} \quad (\text{for}\ 0\leq \alpha\leq k-1), \qquad J \pa_{x^{A, (k)}} = 0,
\end{gather*}
so the image of $J^i=\underbrace{J\circ \cdots \circ J}_{i\text{-times}}$ coincides with the kernel of $J^{k+1-i}$ for $i=1, \ldots, k$. Thus we have an exact sequence
\begin{gather*}
0\rightarrow \operatorname{im} J^k \rightarrow \sT \sT^{k} M \xrightarrow{J} \sV \sT^k M.
\end{gather*}
The image of $J^k$ is given in $\sT \sT^{k}M$ by vanishing all dif\/ferentials $d x^{A, (\alpha)}$ with $0\leq \alpha <k$, hence $\sT \tau^k_{k-1} \colon \sT \sT^k M \twoheadrightarrow \sT\sT^{k-1} M$ vanishes on the kernel of $J$, and so it factors to a vector bundle morphism $\sV \sT^k M \twoheadrightarrow \sT \sT^{k-1} M$. It is immediate to check that the obtained projection coincides with $p^{\sV F_k}_{\pLinr F_k}\colon \sV F_k \rightarrow \pLinr(F_k)$, where $F_k=\sT^k M$.
\end{proof}

The linearisation functor associates in a canonical functorial way a $\GrL$-bundle with any graded bundle. It is clear that one can now iterate the application of the linearisation functor and from a graded bundle of degree $k$ produce a $k$-fold vector bundle. We describe this construction carefully in the next section.

\section[Graded bundles vs $k$-fold vector bundles]{Graded bundles vs $\boldsymbol{k}$-fold vector bundles}\label{sec:GradedBundlesandHigherVectorBundles}
\subsection{The full linearisation functor}

Recall, that the linearisation functor $\pLinr$ takes a graded bundle $F_k$ of degree $k$ to a double graded bundle of bi-degree $(k-1,1)$. We may apply the functor $\pLinr$ again to the graded bundle structure of degree $k-1$ on $\GrL$-bundle $\pLinr(F_k)$ to get a triple graded bundle of multi-degree $(k-2,1,1)$. After $k-1$ iterations we arrive at a $k$-fold vector bundle denoted by $\pLinr^{(k-1)}(F_k)$ which we take as a~def\/inition of the \emph{full linearisation functor}.
\begin{Definition} \label{def:full_linearisation_functor} The full linearisation functor
\begin{gather*}
\Linr\colon \ \GrB[k] \rightarrow \VB[k]
\end{gather*}
is def\/ined as the composition of the linearisation functor $(k-1)$ times, i.e., $\Linr(F_k) = \pLinr^{(k-1)} (F_k)$.
\end{Definition}

\begin{Example}
$\Linr(\sT^k M) = \sT^{(k)} M$ by Proposition~\ref{prop:pLin}.
\end{Example}

Similarly, we can consider iterations of the vertical functor. Recall, the vertical functor $\sV$ takes a graded bundle of degree $k$ to a double graded bundle of degree $(k, 1)$. Applying again the functor $\sV$ to $\big(\sV(F_k), \Delta^1_{\sV F_k}\big)$ we obtain $\sV \sV F_k$ which is a $3$-fold graded bundle of degree $(k, 1, 1)$. Continuing this way we f\/ind that for any $r\geq 1$, $\sV^{(r)} F_k =
\underbrace{\sV \sV\cdots \sV}_{r\text{-times}} F_k$ is a $(r+1)$-fold graded bundle of degree $(k, 1, \ldots, 1)$.

\begin{Proposition}\label{prop:l_r_epi_and_inclusion} There is a canonical epimorphism $(0\leq r\leq k-1)$
\begin{gather}\label{prop:epimorphism_r}
p^{\sV^{(r)}F_k}_{\pLinr^{(r)}({F}_k)}\colon \ \sV^{(r)} F_k \rightarrow \pLinr^{(r)} (F_k) = \sV^{(r)} F_k \big[\Delta^1_{\sV^{(r)} F_k} \leq k-r\big]
\end{gather}
and an embedding
\begin{gather}\label{prop:embedding_r}
\pLinr^{(r)} (F_k) \hookrightarrow \sT^{(r)}\sT^{k-r} F_k.
\end{gather}
Both mappings respect $(r+1)$-fold graded bundle structures.
\end{Proposition}
\begin{proof} We shall prove \eqref{prop:epimorphism_r} by induction. The case $r=0$ is trivial and the case $r=1$ follows directly from the def\/inition of the functor $\pLinr$. Denote $\tilde{F}_k := \sV^{(r)} F_k$ which we consider now as a~graded bundle of degree $k$. By the inductive hypothesis, the base of the canonical projection $\tilde{\tau}^k_{k-r}\colon \tilde{F}_k\rightarrow \tilde{F}_{k}\big[\Delta^1_{\tilde{F}_k \leq k-r}\big]$, which is def\/ined in~\eqref{eqn:fibrations}, coincides with $\pLinr^{(r)}(F_k)$. To get $\pLinr^{(r+1)} F_k$, apply the functor $\pLinr$ to the projec\-tion~$\tilde{\tau}^k_{k-r}$, then compose it with the canonical projection~$p^{\sV \tilde{F}_k}_{\pLinr(\tilde{F}_k)}$ on the linearisation of $\tilde{F}_k$. Using Proposition~\ref{prop:pLin}(\ref{prop_pL:item_pl_tau}) we arrive at the sequence
\begin{gather*}
\sV^{(r+1)} F_k = \sV \tilde{F}_k \xrightarrow{p^{\sV \tilde{F}_k}_{\pLinr(\tilde{F}_k)}} (\sV \tilde{F}_k) \big[\Delta^1_{\sV^{(r+1)}F_k}\leq k-1\big] = \pLinr(\tilde{F}_k) \xrightarrow{\pLinr(\tilde{\tau}^k_{k-r})} \pLinr^{(r+1)} (F_k)  \\
\hphantom{\sV^{(r+1)} F_k = \sV \tilde{F}_k \xrightarrow{p^{\sV \tilde{F}_k}_{\pLinr(\tilde{F}_k)}}}{}
\overset{(*)}{=} \pLinr (\tilde{F}_k) \big[\Delta^1_{\pLinr (\tilde{F}_k)}\leq k-r-1\big] = \sV^{(r+1)} F_k \big[\Delta^1_{\sV^{(r+1)} F_k}\leq k-r-1\big]
\end{gather*}
that completes our inductive reasoning.

To prove~\eqref{prop:embedding_r} consider the graded bundle $\tau\colon F_k\rightarrow M$ as a canonical substructure of $\tau^k_N\colon \sT^k N\rightarrow N$ with $N=F_k$ (for this fundamental observation see~\cite{JG_MR_higher_vb})  and apply the functor $\pLinr$ to both structures. By Proposition~\ref{prop:pLin}\eqref{prop_pL:itemTkM} and~\eqref{prop_pL:item_injective} we get a $\GrL$-bundle embedding $\pLinr(\iota)\colon \pLinr(F_k) \rightarrow \pLinr(\sT^k N) = \sT \sT^{k-1} N$, where $\iota\colon F_k\rightarrow \sT^k N$ denotes the canonical embedding. We can iterate the application of the linearisation functor and f\/inally get~\eqref{prop:embedding_r}.
\end{proof}

\begin{Example}
Consider a graded bundle $F_{2}$ equipped with homogeneous local coordinates $(x^{A}, ~ y^{a}, ~ z^{i})$ of weight $0,1$ and $2$, respectively. Then, $\sV F_2\subset \sT F_2$, $\sV^{(2)}F_{2} \subset \sT^{(2)} F_2$, and $\sV^{(3)} F_2\subset \sT^{(3)} F_2$ come equipped with homogeneous coordinates
\begin{gather*}
\big(x^{A,(0)}_0, y^{a,(0))}_1, z^{i,(0)}_2; y^{a,(1)}_0, z^{i,(1)}_1\big), \\
\big(x^{A, (0,0)}_{0}, y^{a, (0,0)}_{1}, z^{i, (0,0)}_{2}, y^{a, (0,1)}_{0}, z^{i, (0,1)}_{1}; y^{a,(1, 0)}_{0}, z^{i, (1,0)}_{1}, z^{i, (1,1)}_{0}\big), \\
\big(x^{A, (000)}_{0}, y^{a, (000)}_1, z^{i, (000)}_2, y^{a, (001)}_0, z^{i, (001)}_1, y^{a,(010)}_0, z^{i, (010)}_1, z^{i, (011)}_0;  \\
\qquad {} y^{a,(100)}_0, z^{i,(100)}_1, z^{i,(101)}_0, z^{i,(110)}_0\big),
\end{gather*}
respectively. The f\/irst component written in the bottom of the multi-weight is that naturally inherited from $F_{2}$, while the other components come from the vector bundle structure of the tangent bundles. In full generality, consider a graded bundle $F_{k}$ equipped with homogeneous local coordinates $\big(x^{A}, y^i_w\big)$, then $\sV^{(r)}F_{k}$ can be equipped with inherited from~$\sT^{(r)} F_k$ homogeneous local coordinates
\begin{gather*}
\big(x^A, y^{i, (\ep)}_w\big),
\end{gather*}
where $\veps\in\{0,1\}^{r}$ and $|\ep|\leq w$. The weight of $y^{i, (\ep)}_w$ with respect to $\Delta^1_{\sV^{(r)}} F_k$ is $w - |\ep|$ (not explicitly written now). Consider $(\sT^{(r)} F_k; \Delta^0, \Delta^1, \ldots, \Delta^r)$ as a $(r+1)$-fold graded bundle with Euler vector f\/ields $\Delta^1, \ldots, \Delta^r$ of $r$-th iterated tangent bundle of $F_k$ (namely, $\Delta^i$ is the Euler vector f\/ield of the vector bundle projection $\sT^{(i-1)} \tau_{T^{(r-i)} F_k}\colon \sT^{(r)} F_k \rightarrow \sT^{(r-1)} F_k$, $1\leq i\leq r$) and the weight vector f\/ield $\Delta^0:= \dd_{\sT}^{(r)}\Delta_{F_k}$ being the (iterated) tangent lift of the weight vector f\/ield~$\Delta_{F_k}$ of~$F_k$. Def\/ine~$X_r \in \mathfrak{X}(\sT^{(r)} F_k)$ as $X_r = \Delta^0 - \sum_1^k \Delta^k$, so the weight of $y^{a, (\ep_1, \ldots, \ep_r)}_w$ with respect to $X_r$ is $w-(\ep_1+\cdots+\ep_r)$, hence
 \begin{gather}\label{eqn:V_r}
 \sV^{(r)} F_k = \sT^{(r)}F_k[X_r\geq 0].
 \end{gather}
Note that $\Delta^{1}_{\sV^{(r)}F_{k}} = X_{r}|_{\sV^{(r)}F_{k}}$.
\end{Example}

\begin{Example}\label{exmpl:F3}
Let us consider in some detail the generic case when we are dealing with a graded bundle of order three. This example should be enough to highlight the main constructions explicitly without too much notational clutter. Let us employ homogeneous local coordinates $\big(x^{A}, y^{a}, z^{i},w^{\kappa}\big)$ on $F_{3}$ where the weight is assigned as $0$, $1$, $2$ and $3$, respectively. The admissible changes of local coordinates are of the form
\begin{gather*}
 x^{A'} = x^{A'}(x), \qquad  y^{a'} = y^{b}T_{b}^{\:\:\: a'}(x), \qquad  z^{i'} = z^{j}T_{j}^{\:\:\: i'}(x) + \frac{1}{2}y^{a}y^{b}T_{ba}^{\:\:\:\: i'}(x),\\
w^{\kappa'} = w^{\mu}T_{\mu}^{\:\: \kappa'}(x) + z^{i}y^{a}T_{ai}^{\:\:\: \kappa'}(x) + \frac{1}{3!}y^{a}y^{b}y^{c}T_{cba}^{\:\:\:\:\:\: \kappa'}(x). \end{gather*}

Then on $\pLinr(F^3)$ we can employ coordinates $(x^A, \dot{y}^a, y^a, \dot{z}^i,z^i, \dot{w}^\kappa)$ of weights $(0,0)$, $(0,1)$, $(1, 0)$, $(1, 1)$, $(2,0)$ and $(2, 1)$, respectively. The additional transformation rules are
\begin{gather*}
\dot{y}^{a'} = \dot{y}^{b} T^{a'}_{b}, \qquad \dot{z}^{i'}=\dot{z}^j T^{i'}_j + \dot{y}^a y^b T^{i'}_{ba}, \\
\dot{w}^{\kappa'} = \dot{w}^{\mu}T_{\mu}^{\:\: \kappa'} + (\dot{z}^{i}y^{a} + z^{i}\dot{y}^{a}) T_{ai}^{\:\:\: \kappa'} + \frac{1}{2!} \dot{y}^{a}y^{b}y^{c}T_{cba}^{\:\:\:\:\:\: \kappa'}.
\end{gather*}
 By iterating this construction and dif\/ferentiation above formulas, we get $\pLinr^{(2)} F_3$ with coordinates
 $(x^A$, $\dot{y}^a=y^{a,(10)}$, $y^a=y^{a, (00)}$, $\dot{z}^i=z^{i,(10)}$, $dy^a=y^{a,(01)}$, $d\dot{z}^i=z^{i,(11)}$, $dz^i=z^{i,(01)}$, $d\dot{w}^\kappa=w^{\kappa,(11)})$ of degrees $(000)$, $(010)$, $(100)$, $(110)$, $(001)$, $(011)$, $(101)$, $(111)$, respectively, and transformation laws for extra coordinates
\begin{gather*}\allowdisplaybreaks
dy^{a'} = dy^b T^{a'}_b(x), \\
d\dot{z}^{i'} = d\dot{z}^j  T^{i'}_j(x) + \dot{y}^a dy^b T^{i'}_{ba}(x),\\
d z^{i'} = d z^j T^{i'}_j(x) + dy^a y^b T^{i'}_{ba}(x),\\
d \dot{w}^{\kappa'} = d \dot{w}^\mu\,T^{\kappa'}_\mu(x) + (\dot{z}^i dy^a + d\dot{z}^i y^a + dz^i \dot{y}^a)T^{\kappa'}_{ai}(x) +
\dot{y}^a\,dy^b y^c T^{\kappa'}_{cba}(x).
\end{gather*}
We see that transformation laws for coordinates of the same total degree are essentially the same. We conclude that $\pLinr^{(2)}(F_3)$ is a triple vector bundle with isomorphic feet $\pLinr(F_2)$
\begin{gather*}
\xymatrix{
& \pLinr F_2 \ar[rr]\ar[dd] && F_1\ar[dd] \\
\Linr F_3 = \pLinr^{(2)} F_3 \ar[dd] \ar[ru]\ar[rr] && \pLinr F_2 \ar[ru]\ar[dd] & \\
& F_1 \ar[rr]&& M \\
\pLinr F_2 \ar[ru] \ar[rr]&& F_1 \ar[ru] &}
\end{gather*}
\end{Example}

In full generality, it is natural to use the coordinate system $\big(x^A, y^{a, (\ep)}_w\big)$, $\ep\in\{0,1\}^r$, for $\pLinr^{(r)} F_k$ as well (now we take only such $y^{a, (\ep)}_w$ that $0\leq w-|\ep| \leq k-r$), as these coordinates are constant on f\/ibers of the f\/ibration $p^{\sV^{(r)}(F_k)}_{\pLinr^{(r)}(F_k)}$.

On the other hand, consider $F_3$ as a substructure of $\sT^3 F_3$ given by equations $x^{A, (\beta)} = 0$, $y^{a, (\alpha)}_w = 0$, for $1\leq \alpha, \beta\leq 3$ such that $\alpha \neq w$. Then we recognise $\pLinr(F_3)$ as a subset of $\sT \sT^2 F_3$ and $\Linr(F_3)=\pLinr^{(2)}(F_3)$ as a subset of $\sT^{(3)} F_3$, the latter is def\/ined by equations $x^{A, (\ep)} =0$, for $\ep\neq (0,0,0)$ and $y_w^{a, (\ep)} = 0$ if $\ep_1+\ep_2+\ep_3\neq w$ in local coordinate system $\big(x^A, y_w^{a, (\ep)}\big)$, $\ep = (\ep_1, \ep_2, \ep_3)\in\{0, 1\}^3$, on $\sT^{(3)} F_3$. Transformations for the remaining coordinates are the same as given above for $\pLinr^{(2)} F_3$ with the identif\/ication $\big(x^A, y_{|\ep|}^{(\ep_1, \ep_2, \ep_3)}\big) \mapsto \big(x^A, y_{|\ep|}^{(\ep_2, \ep_3)}\big)$, $|\ep| = \ep_1+\ep_2+\ep_3$. The weight of $y_w^{a, (\ep)}$ with respect to the vector f\/ield $X_3\in \mathfrak{X}(F_3)$, def\/ined as the dif\/ference between $\dd_\sT^{(3)} \Delta_{F_3}$ and the sum of three Euler vector f\/ields def\/ining triple vector bundle structure on $\sT^{(3)} F_3$, is equal $w-(\ep_1+\ep_2+\ep_3)$. Thus $\Linr(F_3)$ as a subset of $\sT^{(3)} F_3$ is the set of singular points of the vector f\/ield~$X_3$.

In general, the embedding $\Linr (F_k) \subset \sT^{(k)} F_k$ def\/ined in~\eqref{prop:embedding_r} can be given an explicit form, namely
\begin{gather}\label{eqn:full_linearisation_def}
\Linr(F_k) \simeq \sT^{(k)}F_k[X_k=0],
\end{gather}
where $X_k = \Delta^0 - \sum_1^k \Delta^k$ is as above. (It is not worth to consider $\Linr(F_k)$ as a $(k+1)$-fold graded bundle since $\Delta^0$ restricted to $\Linr(F_k)$ coincides with the sum of Euler vector f\/ields $\Delta^i$, so $(\sT^{(k)} F_k [X = 0]; \Delta^{0})$ is just a graded bundle of degree $k$ def\/ined by taking the total degree in $k$-fold vector bundle~$\Linr(F_k)$.)

Similarly, the epimorphism $\sV^{(k-1)} F_k \rightarrow \Linr(F_k)$ has the explicit form
\begin{gather}\label{eqn:LF_k_explicit}
\sV^{(k-1)} F_k \rightarrow \big(\sV^{(k-1)}F_{k} \big)\big[\Delta^{1}_{\sV^{(k-1)}F_{k}} \leq 1\big]
\end{gather}
and $\Delta^{1}_{\sV^{(k-1)}F_{k}} = X_{k-1}|_{\sV^{(k-1)}F_{k}}$. The isomorphism $\big(V^{(k-1)}F_{k} \big)[\Delta^{1}_{V^{(k-1)}F_{k}} \leq 1] \simeq \sT^{(k)}F_k[X=0]$ has the form
\begin{gather}\label{eqn:local_iso_LFk_explicit}
\big(x^A, y^{i, (\ep)}_w\big) \mapsto \big(x^A, y^{i, (w-|\ep|, \ep_2, \ldots, \ep_k)}_w\big),
\end{gather}
where $\ep = (\ep_2, \ldots, \ep_k)\in \{0,1\}^{k-1}$, so $|\ep| = \ep_2+ \cdots+ \ep_k$, and $0\leq w-|\ep|\leq 1$.

\begin{Corollary}\label{cor:L_emb_and_epim} The full linearisation of a graded bundle $F_k$ can be seen as a quotient structure of the iterated vertical bundle $\sV^{(k-1)} F_k$ and as a substructure of the iterated tangent bundle~$\sT^{(k)} F_k$, as well. The local coordinate formulas are obtained from~\eqref{eqn:V_r}, \eqref{eqn:full_linearisation_def}, \eqref{eqn:LF_k_explicit} and~\eqref{eqn:local_iso_LFk_explicit}.
\end{Corollary}

In view of Corollary~\ref{cor:L_emb_and_epim}, it does not fundamentally matter if we chose to interpret the induced homogeneous coordinates on $\Linr(F_{k})$ as coming from $\sV^{(k-1)}F_{k}$ or $\sT^{(k)}F_{k}$, though the latter maybe easier to comprehend. However, when examining the structure of the full linearisation, understanding~$\Linr(F_{k})$ as a substructure of~$\sT^{(k)}F_{k}$ will prove useful.

